% ASSERT_CONVENTION: natural_units=natural, metric_signature=mostly_minus, fourier_convention=physics, coupling_convention=alpha_s, generator_normalization=T(fund)=1/2, covariant_derivative_sign=D_mu=partial_mu-igA
%
% Lecture 6: Non-SUSY Duality Extensions
% Part of: The Dualised Standard Model -- Part III Lecture Notes
%
% Conventions (Seiberg hep-th/9411149, unchanged from Lectures 1-5):
%   T(fund) = 1/2 per Weyl fermion
%   b_0 = (11N_c - 2N_f)/3 for N_f Dirac flavours
%   A(fund) = 1 (cubic anomaly coefficient per fundamental Weyl fermion)
%   alpha_s = g_s^2 / (4 pi)
%   R[Q] = (N_f - N_c)/N_f; R[W] = 2
%   N_f counts Dirac flavour pairs (each = 2 Weyl fermions)
%
% EPISTEMIC STATUS: This lecture presents CONJECTURAL non-SUSY
% duality extensions.  Every non-SUSY duality claim carries an
% explicit conjectural qualifier ("if the duality holds", etc.).
% This is contract requirement P1 -- NON-NEGOTIABLE.
%
\documentclass[11pt,a4paper]{article}
% ASSERT_CONVENTION: natural_units=natural, metric_signature=mostly_minus, fourier_convention=physics, coupling_convention=alpha_s, generator_normalization=T(fund)=1/2, covariant_derivative_sign=D_mu=partial_mu-igA
%
% CONVENTIONS (Seiberg hep-th/9411149)
% Metric: (+,-,-,-) (mostly minus)
% T(fund) = 1/2 per Weyl fermion
% b_0 = (11N_c - 2N_f)/3 for N_f Dirac flavours
% D_mu = partial_mu - ig A_mu^a T^a
% A(fund) = 1 (cubic anomaly coefficient per fundamental Weyl fermion)
% alpha_s = g_s^2 / (4 pi)
% R[Q] = (N_f - N_c)/N_f; R[W] = 2
% N_f counts Dirac flavour pairs (each = 2 Weyl fermions)
%
% This file is \input{} by all lecture files.

%% ---------- Document class ----------
% (loaded by the wrapper; preamble only provides packages and macros)

%% ---------- Standard packages ----------
\usepackage{amsmath,amssymb,amsthm,mathtools}
\usepackage{hyperref}
\usepackage[capitalise,noabbrev]{cleveref}
\usepackage{enumitem}
\usepackage{booktabs}
\usepackage{array}
\usepackage{xcolor}
\usepackage{tikz}
\usetikzlibrary{arrows.meta,decorations.markings,positioning,calc}
\usepackage{fancybox}
\usepackage{tcolorbox}
\tcbuselibrary{skins,breakable}
\usepackage{slashed}              % Feynman slash notation
\usepackage{cite}                 % compressed citation lists

%% ---------- Hyperref configuration ----------
\hypersetup{
  colorlinks=true,
  linkcolor=blue!70!black,
  citecolor=green!50!black,
  urlcolor=blue!80!black,
  pdfauthor={Part III Lecture Notes},
  pdftitle={The Dualised Standard Model}
}

%% ---------- Theorem environments ----------
\theoremstyle{plain}
\newtheorem{theorem}{Theorem}[section]
\newtheorem{lemma}[theorem]{Lemma}
\newtheorem{proposition}[theorem]{Proposition}
\newtheorem{corollary}[theorem]{Corollary}

\theoremstyle{definition}
\newtheorem{definition}[theorem]{Definition}
\newtheorem{example}[theorem]{Example}
\newtheorem{exercise}{Exercise}[section]

\theoremstyle{remark}
\newtheorem{remark}[theorem]{Remark}

%% ---------- Physics macros: gauge groups and representations ----------
\newcommand{\Nc}{N_{\!c}}
\newcommand{\Nf}{N_{\!f}}
\newcommand{\SU}[1]{\mathrm{SU}(#1)}
\newcommand{\UU}[1]{\mathrm{U}(#1)}

% Representations
\newcommand{\fund}{\square}
\newcommand{\antifund}{\overline{\square}}
\newcommand{\adj}{\mathrm{adj}}
\newcommand{\antisym}{\wedge^{\!2}}        % antisymmetric rank-2
\newcommand{\sym}{\mathrm{S}_2}           % symmetric rank-2

%% ---------- Physics macros: group theory constants ----------
% Dynkin index: Tr(T^a T^b) = T(R) delta^{ab}
\newcommand{\Tfund}{\tfrac{1}{2}}          % T(fund) = 1/2 per Weyl fermion
\newcommand{\Tadj}{\Nc}                    % T(adj) = N_c
\newcommand{\Cfund}{\frac{\Nc^2 - 1}{2\Nc}}  % C_2(fund)
\newcommand{\Cadj}{\Nc}                    % C_2(adj)

% Cubic anomaly coefficient: A(R) = Tr_R[T^a {T^b, T^c}] with A(fund) = 1
\newcommand{\Afund}{1}

%% ---------- Physics macros: beta function ----------
\newcommand{\bzero}{b_0}
\newcommand{\bone}{b_1}
\newcommand{\alphas}{\alpha_s}
\newcommand{\gs}{g_s}

%% ---------- Physics macros: SUSY / R-charge ----------
\newcommand{\Rcharge}[1]{R[#1]}
\newcommand{\Wsup}{W}                      % superpotential

%% ---------- Physics macros: covariant derivative and field strength ----------
\newcommand{\Dmu}{D_\mu}
\newcommand{\Fmn}{F_{\mu\nu}}
\newcommand{\Ftilde}{\widetilde{F}_{\mu\nu}}

%% ---------- Physics macros: common abbreviations ----------
\newcommand{\ie}{\textit{i.e.}}
\newcommand{\eg}{\textit{e.g.}}
\newcommand{\cf}{\textit{cf.}}
\newcommand{\IR}{\ensuremath{\mathrm{IR}}}
\newcommand{\UV}{\ensuremath{\mathrm{UV}}}
\newcommand{\AF}{\ensuremath{\mathrm{AF}}}
\newcommand{\BZ}{\ensuremath{\mathrm{BZ}}}
\newcommand{\NSVZ}{\ensuremath{\mathrm{NSVZ}}}
\newcommand{\SQCD}{\ensuremath{\mathrm{SQCD}}}
\newcommand{\QCD}{\ensuremath{\mathrm{QCD}}}
\newcommand{\SM}{\ensuremath{\mathrm{SM}}}
\newcommand{\Tr}{\mathrm{Tr}}

%% ---------- Convention box macro ----------
\newtcolorbox{conventionbox}[1][]{%
  enhanced,
  colback=blue!5!white,
  colframe=blue!60!black,
  fonttitle=\bfseries,
  title={Conventions},
  #1
}

%% ---------- Comparison table style ----------
\newtcolorbox{comparisontable}[1][]{%
  enhanced,
  colback=orange!5!white,
  colframe=orange!70!black,
  fonttitle=\bfseries,
  title={Comparison},
  #1
}

%% ---------- Warning / important box ----------
\newtcolorbox{warningbox}[1][]{%
  enhanced,
  colback=red!5!white,
  colframe=red!70!black,
  fonttitle=\bfseries,
  title={Important},
  #1
}

%% ---------- Preview / forward-reference box ----------
\newtcolorbox{previewbox}[1][]{%
  enhanced,
  colback=green!5!white,
  colframe=green!50!black,
  fonttitle=\bfseries,
  title={Preview},
  #1
}

%% ---------- Page layout ----------
\usepackage[margin=2.5cm]{geometry}
\setlength{\parskip}{0.4em}
\setlength{\parindent}{0em}

%% ---------- Section formatting ----------
\usepackage{titlesec}
\titleformat{\section}{\Large\bfseries}{\thesection}{1em}{}
\titleformat{\subsection}{\large\bfseries}{\thesubsection}{1em}{}

%% ---------- End of preamble ----------


\title{\textbf{Lecture 6: Non-SUSY Duality Extensions}}
\author{The Dualised Standard Model --- Part III Lecture Notes}
\date{Lent Term 2027}

\begin{document}
\maketitle

\tableofcontents
\newpage

%% ========================================================================
%%  CONVENTIONS
%% ========================================================================

\begin{conventionbox}[title={Conventions for Lecture 6}]
All conventions from Lectures~1--5 remain in force.  No new conventions
are introduced in this lecture.

\renewcommand{\arraystretch}{1.3}
\begin{tabular}{@{}l l l@{}}
\toprule
\textbf{Quantity} & \textbf{Convention} & \textbf{Comment} \\
\midrule
Units & $\hbar = c = 1$ (natural) & \\
Metric & $\eta_{\mu\nu} = \mathrm{diag}(+1,-1,-1,-1)$ & mostly minus \\
Dynkin index & $\Tr(T^a T^b) = \Tfund\, \delta^{ab}$ & fundamental of $\SU{\Nc}$ \\
Cubic anomaly & $A(\fund) = 1$, $A(\antifund) = -1$ & per fundamental Weyl \\
Beta function & $\bzero = (11\Nc - 2\Nf)/3$ & $\Nf$ Dirac flavours \\
R-charge & $\Rcharge{Q} = (\Nf - \Nc)/\Nf$ & SUSY reference point \\
$\Nf$ counting & Dirac flavour pairs & each $= 2$ Weyl fermions \\
\bottomrule
\end{tabular}

\medskip

\textbf{Epistemic convention for this lecture.}  Every statement
asserting or assuming non-supersymmetric duality carries an explicit
conjectural qualifier: ``if the duality holds,'' ``assuming the
conjectured duality,'' ``conjectural,'' or equivalent.  The anomaly
matching \emph{equations} presented below are exact; the
\emph{duality interpretation} of those equations is conjectural.
\end{conventionbox}


%% ========================================================================
%%  SECTION 1: What We Lose Without SUSY
%% ========================================================================
\section{What We Lose Without SUSY}
\label{sec:what-we-lose}

In Lectures~2--4 we constructed Seiberg duality for $\mathcal{N}=1$
SQCD as an equivalence established through exact consistency checks.  That
case rested on four independent pillars working in concert.  Before
extending duality beyond supersymmetry, we must be explicit about
which of these pillars survive and which do not.

\subsection{The four pillars of SUSY duality}

\begin{enumerate}[label=(\roman*),leftmargin=*]
  \item \textbf{Holomorphy of the superpotential} (Lecture~2).  The
    superpotential $\Wsup$ is a holomorphic function of the chiral
    superfields and is not renormalised beyond one loop.  This
    non-renormalisation theorem~\cite{Seiberg1993} constrains the
    effective superpotential exactly, leaving the dynamical scale
    $\Lambda$ as the only non-perturbative input.

  \item \textbf{NSVZ exact beta function} (Lecture~2).  The exact
    relation between the beta function, the anomalous dimensions, and
    the gauge coupling~\cite{NSVZ1983} determines the boundaries of
    the conformal window \emph{exactly}: $\tfrac{3}{2}\Nc < \Nf < 3\Nc$
    for $\SU{\Nc}$ with $\Nf$ fundamental flavours.

  \item \textbf{Moduli space matching} (Lecture~3).  The classical
    moduli space of flat directions is not lifted quantum mechanically
    (in the conformal window), and the electric and magnetic moduli
    spaces can be mapped one-to-one via the meson operator
    $M^i_j = Q^i \widetilde{Q}_j / \Lambda$.  This provides an
    infinite family of consistency checks---one for each point on the
    moduli space.

  \item \textbf{'t~Hooft anomaly matching} (Lecture~1, verified
    in Lecture~3).  All six independent anomaly
    coefficients---$\SU{\Nf}_L^3$, $\SU{\Nf}_R^3$,
    $\SU{\Nf}_L^2\,\UU{1}_B$, $\SU{\Nf}_R^2\,\UU{1}_B$,
    $\UU{1}_B^3$, $\UU{1}_B$-gravitational---match exactly between
    the electric and magnetic descriptions as polynomial identities in
    $\Nc$ and $\Nf$ (validated by \texttt{verify\_anomaly\_matching.py},
    7/7 symbolic checks).
\end{enumerate}

\subsection{What is lost}

Without supersymmetry, three of the four pillars collapse:

\begin{enumerate}[label=(\alph*)]
  \item \textbf{Holomorphy:} There is no superpotential in a non-SUSY
    theory.  The powerful non-renormalisation theorems that fix the
    form of $\Wsup$ have no analogue.  All couplings run and mix
    under renormalisation without holomorphic constraints.

  \item \textbf{Exact beta function:} The NSVZ formula relies on the
    superfield structure.  Without SUSY, only the perturbative
    coefficients $\bzero$, $\bone$, \ldots\ are available.  The
    boundaries of the conformal window become \emph{estimates} rather
    than exact results.

  \item \textbf{Moduli space:} Flat directions require cancellation of
    the scalar potential, which in SUSY theories is guaranteed by the
    $F$-term and $D$-term conditions.  Without SUSY, quantum
    corrections generically lift all flat directions.  There is no
    moduli space to match.
\end{enumerate}

\subsection{What survives}

\begin{warningbox}[title={The sole surviving tool}]
  Of the four pillars, \textbf{only 't~Hooft anomaly matching
  survives} without SUSY.

  Anomaly matching is a non-perturbative, exact constraint that
  relates the UV and IR spectra of \emph{any} gauge theory.  It was
  derived by 't~Hooft~\cite{tHooft1980} in 1980 without any
  reference to supersymmetry and requires only:
  \begin{enumerate}[label=(\roman*)]
    \item a UV gauge theory with a global symmetry group $G_{\mathrm{global}}$,
    \item massless fermions charged under $G_{\mathrm{global}}$
      in both the UV and IR descriptions.
  \end{enumerate}
  This is exactly the tool we developed in Lecture~1 and applied in
  Lecture~3 to verify Seiberg duality.  It remains available---and
  equally powerful---in the absence of SUSY.
\end{warningbox}

\begin{comparisontable}[title={SUSY vs.\ non-SUSY evidence for duality}]
\renewcommand{\arraystretch}{1.5}
\begin{tabular}{@{} >{\raggedright}p{4cm} >{\centering}p{3cm}
  >{\centering\arraybackslash}p{3cm} @{}}
\toprule
\textbf{Evidence type} & \textbf{SUSY duality} & \textbf{Non-SUSY duality} \\
\midrule
Holomorphy & \checkmark\ (exact) & $\times$ \\
Exact beta function (NSVZ) & \checkmark\ (exact) & $\times$ \\
Moduli space matching & \checkmark\ (infinite checks) & $\times$ \\
't~Hooft anomaly matching & \checkmark\ (exact) & \checkmark\ (exact) \\
Superconformal index & \checkmark\ (independent check) & $\times$ \\
\midrule
\textbf{Total independent pillars} & \textbf{5} & \textbf{1} \\
\bottomrule
\end{tabular}

\medskip

\textbf{Summary:} Non-SUSY duality proposals rest on anomaly matching
alone.  Anomaly matching is \emph{necessary} but \emph{not sufficient}
for duality: multiple distinct IR spectra can satisfy the same UV
anomaly conditions.  Every non-SUSY duality statement in this lecture
is therefore \textbf{conjectural}.
\end{comparisontable}


%% ========================================================================
%%  SECTION 2: Sannino-Schechter Toy Model (1997)
%% ========================================================================
\section{The Sannino--Schechter Toy Model (1997)}
\label{sec:sannino-schechter}

\textit{Status: toy model --- proof of concept, not a derivation of non-SUSY duality.}

The first concrete evidence that SUSY duality structures might survive
the breaking of supersymmetry came from a toy model by Sannino and
Schechter~\cite{SS1997}.

\subsection{The key insight}

Sannino and Schechter started from the effective Lagrangian for
$\mathcal{N}=1$ super Yang--Mills theory (the Veneziano--Yankielowicz
Lagrangian~\cite{VY1982}) and asked: what happens if we add
supersymmetry-breaking terms and integrate out the heavy superpartners
(gluinos and squarks)?

Their key observation was that the ordinary QCD degrees of
freedom---mesons and glueballs---\emph{emerge from the auxiliary fields}
of the supersymmetric effective Lagrangian.  Specifically:
\begin{itemize}
  \item The scalar and pseudoscalar glueball fields of QCD are
    contained in the auxiliary $F$-field of the chiral superfield $S$
    (see Eq.~(2.4) of~\cite{SS1997}):
    \begin{equation}\label{eq:F-field}
      F \;=\; \frac{g^2}{32\pi^2}\Bigl[
        -\tfrac{1}{2} F_a^{mn} F_{mn,a}
        - \tfrac{i}{4}\epsilon_{mnrs} F_a^{mn} F_a^{rs}
        + \cdots
      \Bigr]\,.
    \end{equation}
  \item The meson fields of QCD are hidden in the auxiliary field
    $(F_T)_{ij} \sim \psi_i \widetilde{\psi}_j$ of the meson
    superfield.
\end{itemize}

When SUSY is broken at a scale $m$ and the superpartners are
integrated out, these auxiliary fields become the \emph{dynamical}
degrees of freedom of the low-energy theory.  Sannino and Schechter
showed that the resulting effective Lagrangian, under suitable
assumptions about the SUSY-breaking terms, reproduces features of the
ordinary QCD effective Lagrangian.

\subsection{What this suggests---and what it does not prove}

The Sannino--Schechter model \emph{suggests} that the duality structure
encoded in the SUSY effective Lagrangian may persist, in some form,
after SUSY breaking.  The QCD-like degrees of freedom that emerge are
not ad hoc additions but arise organically from the SUSY structure.
Sannino and Schechter themselves described this as an ``amusing feature''
of their model~\cite{SS1997}.

However, the model does \emph{not} prove that non-SUSY duality exists.
The SUSY-breaking procedure involves arbitrary choices (the form of
the breaking terms, the assumption that the dominant part of the
effective Lagrangian retains a ``tree-level holomorphicity''), and the
resulting QCD Lagrangian depends on these choices.  The model is a
proof of concept, not a derivation.

\subsection{The Cheng--Shadmi tension}
\label{ssec:cheng-shadmi}

Shortly after the Sannino--Schechter proposal, Cheng and
Shadmi~\cite{CS1998} performed a systematic analysis of Seiberg
duality in the presence of soft SUSY-breaking terms.  Their approach
was to couple both the electric and magnetic theories to a common
dynamical SUSY-breaking (DSB) sector and track the induced soft masses
through the duality map.

\begin{warningbox}[title={Cheng--Shadmi tension (unresolved)}]
  Cheng and Shadmi found that in the limit of \textbf{large SUSY
  breaking}---precisely the limit relevant for obtaining non-SUSY
  QCD from broken SQCD---the magnetic scalars generically acquire
  \textbf{negative mass-squared}.  This destabilises the magnetic
  vacuum: the magnetic theory does not have a stable ground state in
  the regime where the electric theory reduces to ordinary QCD.

  \medskip

  Specifically, assuming a canonical K\"ahler potential for the
  magnetic theory (the simplest and most natural choice), they showed
  that the soft scalar masses satisfy $m_{\mathrm{scalar}}^2 < 0$
  for many cases in the conformal window~\cite{CS1998}.  The
  resulting symmetry breaking pattern in the magnetic theory is
  \textbf{incompatible} with the expected chiral symmetry breaking
  pattern of QCD.

  \medskip

  \textbf{This tension is unresolved.}  Sannino's subsequent
  programme (Sections~\ref{sec:qcd-dual}--\ref{sec:emergent-susy})
  \emph{sidesteps} rather than \emph{resolves} the Cheng--Shadmi
  problem by conjecturing non-SUSY duality directly, without deriving
  it from broken SUSY.  The Cheng--Shadmi result remains a cautionary
  example of what can go wrong when SUSY-derived structures are
  extrapolated to the non-SUSY regime.
\end{warningbox}


%% ========================================================================
%%  SECTION 3: The Non-SUSY QCD Dual (Sannino, 2009)
%% ========================================================================
\section{The Non-SUSY QCD Dual (Sannino, 2009)}
\label{sec:qcd-dual}

\textit{Status: anomaly-based candidate --- consistency checks passed, but anomaly matching alone is not sufficient for proof.}

Rather than deriving non-SUSY duality from broken SUSY---an approach
plagued by the Cheng--Shadmi tension---Sannino~\cite{Sannino2009}
proposed in 2009 to \emph{conjecture} non-SUSY duality directly and
check it for consistency via anomaly matching.

\subsection{The electric theory}

The starting point is ordinary QCD-like physics: an $\SU{\Nc}$ gauge
theory (specialised to $\Nc = 3$ for QCD) with $\Nf$ Dirac flavours
of fundamental quarks.  The quantum global symmetry group is:
\begin{equation}\label{eq:qcd-global}
  G_{\mathrm{global}}
  \;=\; \SU{\Nf}_L \times \SU{\Nf}_R \times \UU{1}_V\,,
\end{equation}
where the classical axial $\UU{1}_A$ is broken by the Adler--Bell--Jackiw
anomaly (as discussed in Lecture~1).

\subsection{The conjectured magnetic dual}

Sannino sought a \emph{magnetic} gauge theory with gauge group
$\SU{X}$---where $X$ is to be determined by anomaly matching---and a
matter content of magnetic quarks $q$, $\widetilde{q}$ plus
gauge-singlet fermions identifiable as baryons built from the electric
quarks.  The key constraint is that the magnetic spectrum must
reproduce \emph{all} independent 't~Hooft anomaly coefficients
of the electric theory:
\begin{equation}\label{eq:anomaly-conditions}
  \SU{\Nf}_{L/R}^3\,,
  \qquad
  \SU{\Nf}_{L/R}^2\,\UU{1}_V\,.
\end{equation}
(For a vector-like theory, these are the independent anomaly types;
\cf\ Lecture~1.)

\subsection{A non-trivial solution}

Sannino found a novel solution to the anomaly matching conditions
with~\cite{Sannino2009}:
\begin{equation}\label{eq:dual-gauge-group}
  \boxed{X \;=\; 2\Nf - 5N\,,}
\end{equation}
where $N = \Nc = 3$ for QCD, giving a dual gauge group
$\SU{2\Nf - 15}$.  (For $\Nf = 9$ and $N = 3$, the magnetic gauge
group is $\SU{3}$---the same as the electric theory, though with
different matter content and interactions.)

The solution includes magnetic quarks in the fundamental of
$\SU{X}$ plus gauge-singlet Weyl fermions in specific representations
of the global symmetry group, with multiplicities and $\UU{1}_V$
charges determined by the anomaly matching equations.  The full
spectrum is tabulated in Table~III of~\cite{Sannino2009}.

\subsection{Consistency checks}

\emph{If the conjectured duality holds}, Sannino's solution passes
several non-trivial consistency checks:

\begin{enumerate}[label=(\alph*)]
  \item \textbf{Anomaly matching:} All independent 't~Hooft anomaly
    coefficients match between the electric and magnetic descriptions.
    The anomaly matching equations are polynomials in $\Nf$ and $N$,
    and the solution satisfies them as identities.

  \item \textbf{Conformal window prediction:}  Requiring asymptotic
    freedom of the dual theory gives a lower bound on $\Nf$:
    \begin{equation}\label{eq:dual-AF}
      \Nf \;\geq\; \frac{11}{4}\,N\,,
    \end{equation}
    which for $N = 3$ gives $\Nf \geq 8.25$.  Remarkably, this
    \emph{coincides} with the lower boundary of the conformal window
    predicted by the all-orders conjectured beta function of
    Ryttov and Sannino~\cite{RS2008} when the anomalous dimension of
    the mass operator takes the value $\gamma = 2$.

  \item \textbf{Involution:}  Dualising the magnetic theory recovers
    the electric gauge group structure---the duality is its own
    inverse~\cite{Sannino2009}.
\end{enumerate}

\begin{warningbox}[title={Anomaly matching is necessary but not sufficient}]
  The existence of a consistent anomaly matching solution is
  \textbf{evidence for} duality, not \textbf{proof of} it.  The
  anomaly matching conditions are a \emph{finite} set of polynomial
  equations in the multiplicities and quantum numbers of the magnetic
  spectrum.  Multiple distinct solutions can satisfy the same
  conditions~\cite{Sannino2009}.

  \medskip

  Sannino's solution is distinguished by additional
  physics requirements (involution, flavour decoupling, consistency
  with the conjectured beta function), but these are
  \emph{consistency checks}, not derivations.  The duality itself
  remains a conjecture.
\end{warningbox}


%% ========================================================================
%%  SECTION 4: The Standard Model as a Magnetic Dual (Sannino, 2011)
%% ========================================================================
\section{The Standard Model as a Magnetic Dual (Sannino, 2011)}
\label{sec:sm-magnetic}

\textit{Status: conjectural framework --- anomaly matching plus structural consistency, but no independent cross-validation.}

In 2011, Sannino~\cite{Sannino2011} extended the QCD dual framework
to the full Standard Model by embedding it in a Pati--Salam-like
structure.  \emph{If the non-SUSY duality holds}, the Standard Model
can be viewed as the \emph{magnetic dual} of a purely fermionic
electric theory.

\subsection{The SUSY precursor: Maekawa--Takahashi}

\textit{Status: established SUSY construction --- Seiberg duality in the Pati--Salam sector is proven within $\mathcal{N}=1$ SQCD.}

The direct SUSY precursor of this construction was provided by
Maekawa and Takahashi~\cite{MT1996}, who constructed the Seiberg dual
of a supersymmetric theory with Pati--Salam gauge group
$\SU{4}_{\mathrm{PS}}$.  In the SUSY context, their construction
naturally produces:
\begin{itemize}
  \item A top quark Yukawa coupling of order unity (from the magnetic
    superpotential $\Wsup_{\mathrm{mag}} \sim M\,q\,\widetilde{q}$),
    explaining the large top mass without fine-tuning.
  \item A natural $\mu$-parameter for the Higgs sector.
\end{itemize}
Sannino's 2011 proposal extends the Maekawa--Takahashi construction
to non-SUSY by replacing the SUSY consistency arguments with anomaly
matching.

\subsection{The Pati--Salam embedding}

The Standard Model fermion content---quarks and leptons---is naturally
embedded in a Pati--Salam $\SU{4}$ framework by treating the lepton
number as a fourth colour~\cite{PS1974}.  In Sannino's construction,
the \textbf{magnetic} (SM-side) description includes~\cite{Sannino2011}:
\begin{itemize}
  \item Quarks and leptons: $p$ and $\widetilde{p}$ in the
    fundamental and antifundamental of $\SU{4}_{\mathrm{PS}}$,
    carrying flavour indices under the global symmetry group.
  \item The Higgs field $H$: transforming as a bifundamental of the
    flavour group---\emph{not} as an elementary scalar of the UV
    theory, but as a \textbf{composite} built from electric fermions.
  \item An adjoint Weyl fermion $\lambda_m$ of $\SU{4}_{\mathrm{PS}}$.
  \item A meson field $M$: a scalar transforming as a bifundamental
    of the flavour group, mediating interactions between the magnetic
    quarks and the gauge-singlet fermions.
\end{itemize}

\subsection{Three conjectural consequences}

\emph{If the SM is the magnetic dual of a Pati--Salam-like electric
theory}, three remarkable consequences follow:

\begin{enumerate}[label=(\alph*)]
  \item \textbf{No elementary scalars in the UV = no hierarchy
    problem.}  The electric theory contains \emph{only fermions and
    gauge bosons}---no fundamental scalars.  The Higgs doublets of
    the SM are composite ``mesons'' of the electric theory
    ($H \sim Q\lambda\lambda\widetilde{Q}$ schematically), analogous
    to the meson $M$ of Seiberg duality.  Since the electric theory
    has no scalars, there are no quadratic divergences in the
    fundamental description, and the hierarchy problem is, if the
    duality holds, an artefact of the magnetic (composite) description.

  \item \textbf{$N_{\mathrm{gen}} \geq 3$ from the dual gauge group.}
    The dual \emph{electric} gauge group is
    $\SU{2n_g - 4} = \SU{\Nf - 4}$, where $n_g$ is the number of
    SM generations and $\Nf = 2n_g$~\cite{Sannino2011}.  For this
    group to be non-abelian (a necessary condition for a
    non-trivial gauge theory), we need:
    \begin{equation}\label{eq:Ngen}
      2n_g - 4 \;\geq\; 2
      \qquad\Longrightarrow\qquad
      \boxed{n_g \;\geq\; 3\,.}
    \end{equation}
    \emph{If the duality holds}, the existence of at least three
    generations of quarks and leptons is not an accident but a
    \emph{mathematical requirement} for the electric dual to be a
    non-trivial gauge theory.  Requiring asymptotic freedom of the
    magnetic theory further constrains $n_g \leq 6$~\cite{Sannino2011}.

  \item \textbf{Common Yukawa origin.}  The Yukawa couplings of the SM
    arise from the magnetic interaction Lagrangian, which has the
    schematic form $\mathcal{L}_{\mathrm{mag}} \supset
    Y\, q\, M\, \widetilde{q}$---the non-SUSY analogue of the SUSY
    magnetic superpotential $\Wsup_{\mathrm{mag}} = M\,q\,\widetilde{q}/\mu$.
    \emph{If the duality holds}, all Yukawa couplings share a common
    dynamical origin rather than being independent free parameters.
\end{enumerate}

\begin{remark}[Why ``if the duality holds'' matters]
  Each of the three consequences above is conditional on the
  \emph{conjectured} non-SUSY duality being correct.  Without a proof
  of duality, they are \textbf{predictions of the framework}, not
  established results.  They are compelling because they address
  genuine puzzles of the SM (hierarchy problem, number of
  generations, Yukawa structure), but their validity depends entirely
  on whether the underlying duality conjecture is correct.
\end{remark}


%% ========================================================================
%%  SECTION 5: Emergent SUSY at Fixed Points (2011/2013)
%% ========================================================================
\section{Emergent SUSY at Fixed Points (Antipin--Mojaza--Pica--Sannino, 2011)}
\label{sec:emergent-susy}

\textit{Status: large-$N$ perturbative evidence --- established in the Veneziano limit, extrapolation to finite $N$ uncontrolled.}

The most striking piece of perturbative evidence for non-SUSY duality
was provided by Antipin, Mojaza, Pica, and
Sannino~\cite{AMPS2011}, who showed that \emph{supersymmetric fixed
points can emerge from non-supersymmetric RG flow} in the magnetic
dual theory.

\subsection{The setup}

The magnetic dual theory of Section~\ref{sec:qcd-dual}, augmented by
a Weyl adjoint fermion $\lambda_m$ (as in the SM-as-magnetic
construction of Section~\ref{sec:sm-magnetic}), has the quantum global
symmetry:
\begin{equation}\label{eq:global-sym-magnetic}
  \SU{\Nf}_L \times \SU{\Nf}_R \times \UU{1}_V
  \times \UU{1}_{AF}\,,
\end{equation}
where $\UU{1}_{AF}$ is an anomaly-free axial-like symmetry under which
the adjoint fermion is charged.

The theory has multiple couplings: the gauge coupling $\alpha_g$,
Yukawa-type couplings $\alpha_y$ mediating interactions between the
magnetic quarks and the scalar mesons, and quartic scalar couplings.
In the Veneziano limit ($\Nc, \Nf \to \infty$ with $\Nf/\Nc$ fixed),
the perturbative beta functions for all these couplings can be
computed systematically.

\subsection{The result: emergent SUSY}

Antipin \textit{et al.}\ classified all perturbatively accessible
fixed points of the coupled beta function system and discovered that
among a large number of non-supersymmetric fixed points, there exist
fixed points where~\cite{AMPS2011}:
\begin{enumerate}[label=(\roman*)]
  \item the gauge coupling, Yukawa couplings, and quartic scalar
    couplings all take \emph{precisely the values predicted by
    supersymmetry};
  \item these SUSY fixed points are reached by the RG flow starting
    from \emph{generic non-supersymmetric} initial conditions---no
    fine-tuning of bare couplings is required;
  \item the SUSY fixed points have finite basins of attraction in the
    full coupling-constant space.
\end{enumerate}

The perturbative analysis \emph{suggests} that supersymmetry
can emerge dynamically as an IR phenomenon: a non-SUSY UV Lagrangian
flows to a SUSY fixed point without any SUSY being imposed at the
fundamental level.  This supports the ``quasi-supersymmetric''
interpretation of the magnetic spectrum~\cite{Sannino2011}: even
without imposing SUSY, the IR dynamics organises fields into
SUSY-like multiplets at the fixed point.

\subsection{Limitations}

\begin{warningbox}[title={Perturbative result in the Veneziano limit}]
  The emergent SUSY result of Antipin \textit{et al.}\ is a
  \textbf{perturbative} result valid in the \textbf{Veneziano limit}:
  $\Nc, \Nf \to \infty$ with $\Nf/\Nc$ held fixed.  The
  perturbative analysis is reliable because, in this limit, the fixed
  point couplings are parametrically small (controlled by
  $\epsilon = \Nf^{\mathrm{AF}} - \Nf \ll \Nf^{\mathrm{AF}}$).

  \medskip

  Whether emergent SUSY survives at \textbf{finite} $\Nc$ and $\Nf$
  ---and in particular at the SM values $\Nc = 3$,
  $\Nf = 6$ (three generations)---is an \textbf{open question}.
  Finite-$N$ corrections to the beta functions are uncontrolled
  and could qualitatively change the fixed-point structure.

  \medskip

  \textbf{The claim that SUSY emerges at the SM point extrapolates
  a perturbative large-$N$ result to a regime where it is not
  controlled.}  This does not invalidate the result---it means that
  the result is \emph{suggestive}, not \emph{conclusive}.
\end{warningbox}

\subsection{Complementary large-$N$ arguments}

A complementary approach to non-SUSY duality at large $N$ was
provided by Armoni, Shifman, and Unsal~\cite{ASU2008}, who used
\emph{orientifold planar equivalence} to argue that certain non-SUSY
gauge theories share the same planar limit as their SUSY counterparts.
This provides another piece of evidence for non-SUSY duality, but
again limited to the large-$N$ regime.  Whether the planar equivalence
extends to finite $N$ is unknown.


%% ========================================================================
%%  SECTION 6: The Central Vulnerability
%% ========================================================================
\section{The Central Vulnerability}
\label{sec:vulnerability}

We now state, as the concluding point of this lecture, the
\textbf{central vulnerability} of the non-SUSY duality programme.
This is not a footnote or caveat---it is the most important
intellectual content of the lecture.

\subsection{The vulnerability stated}

\begin{warningbox}[title={Central vulnerability of non-SUSY duality}]
  \textbf{Non-SUSY duality has no independent cross-validation
  beyond 't~Hooft anomaly matching.}

  \medskip

  In SUSY, duality is established by \textbf{four independent
  pillars} working together (\S\ref{sec:what-we-lose}):
  anomaly matching, holomorphy, moduli space matching, and the NSVZ
  exact beta function (plus the superconformal index as a fifth check).
  The convergence of multiple independent lines of evidence is what
  makes SUSY duality \emph{proven}.

  \medskip

  In non-SUSY, \textbf{only anomaly matching survives}.  The
  construction \emph{postulates} the duality and checks consistency
  via anomaly matching.  It does not \emph{prove} the duality.

  \medskip

  The Sannino programme is the simplest \emph{consistent} framework
  that:
  \begin{enumerate}[label=(\roman*)]
    \item satisfies all 't~Hooft anomaly conditions,
    \item passes the involution and flavour decoupling tests,
    \item reproduces the correct SM gauge group and matter content,
    \item predicts $n_g \geq 3$ and offers a mechanism for the
      hierarchy problem.
  \end{enumerate}
  But \textbf{consistency is not proof}.  Anomaly matching is
  necessary but not sufficient: multiple IR spectra can satisfy the
  same UV anomaly conditions.
\end{warningbox}

\subsection{Then why study it?}

A student who has followed the rigorous development of SUSY duality
in Lectures~2--4 may reasonably ask: ``If non-SUSY duality is
unproven, why are we studying it?''

The answer involves four considerations:

\begin{enumerate}[label=(\alph*)]
  \item \textbf{Anomaly matching is highly constraining.}  While not
    sufficient by itself, anomaly matching is an extremely powerful
    consistency condition.  The number of solutions that simultaneously
    satisfy all anomaly conditions, pass the involution test, and allow
    flavour decoupling is small.  Finding even one consistent solution
    is non-trivial.

  \item \textbf{The SUSY evidence is suggestive.}  Seiberg duality
    for SQCD is proven.  The physical picture of tumbling plus
    complementarity (Lecture~5) provides a mechanism that does not
    intrinsically require SUSY.  The Sannino--Schechter toy model
    (\S\ref{sec:sannino-schechter}) shows that QCD-like degrees of
    freedom emerge naturally from the SUSY structure.  The emergent
    SUSY result (\S\ref{sec:emergent-susy}) suggests that the
    fixed-point structure survives beyond SUSY.  None of this is proof,
    but collectively it makes the conjecture \emph{well-motivated}.

  \item \textbf{The framework produces non-trivial predictions.}
    \emph{If the duality holds}, it explains the hierarchy problem
    (no fundamental scalars), predicts $n_g \geq 3$ (non-abelian
    electric gauge group), and provides a common origin for Yukawa
    couplings.  These are testable consequences of the framework,
    not tautologies.

  \item \textbf{Belief is not proof.}  We study the framework
    \emph{because} it is the most economical and well-motivated
    conjecture consistent with all available evidence---not because
    it is established.  Independent validation is needed:
    \begin{itemize}
      \item \textbf{Lattice:} Direct numerical verification of the
        non-SUSY conformal window structure and the duality map.
      \item \textbf{New formal tools:} Non-perturbative methods
        beyond anomaly matching that could provide additional
        cross-checks.
      \item \textbf{Experimental signatures:} Predictions of the
        dualised SM (composite Higgs, extra matter at the
        Pati--Salam scale) that could be tested at future colliders.
    \end{itemize}
\end{enumerate}

\subsection{Recent developments}

The non-SUSY duality programme remains an active area of research.
Bi and Terning~\cite{BT2022} have explored non-SUSY duality
candidates for chiral gauge theories with $\mathrm{Spin}(8)$
gauge group, verifying 21 independent anomaly matching conditions.
Their work demonstrates that the programme extends beyond QCD-like
(vector-like) theories to chiral gauge theories---a broader and
more challenging setting.

The backbone paper of Cacciapaglia and
Sannino~\cite{CacciapagliaSannino2024} provides the most
comprehensive treatment of the dualised Standard Model, including the
quasi-supersymmetric spectrum and quantitative IR matching, which we
will develop in Lectures~7--8.


%% ========================================================================
%%  FORWARD REFERENCE
%% ========================================================================

\begin{previewbox}[title={Preview: Lectures 7--8}]
  In Lectures~7--8, we will \emph{assume} the non-SUSY duality
  conjecture and construct the dualised Standard Model following
  Cacciapaglia and Sannino~\cite{CacciapagliaSannino2024}.  We will:
  \begin{itemize}
    \item Specify the electric theory (Pati--Salam gauge group,
      purely fermionic matter content).
    \item Construct the magnetic dual and identify the SM fields as
      composites.
    \item Derive the quasi-supersymmetric spectrum and check anomaly
      matching for the full $\SU{3}_C \times \SU{2}_L \times
      \UU{1}_Y$ gauge group.
    \item Compute the fermion mass relations and compare with
      observation.
  \end{itemize}

  \textbf{Every prediction derived in Lectures~7--8 is conditional on
  the conjectured non-SUSY duality being correct.}
\end{previewbox}


%% ========================================================================
%%  SUMMARY
%% ========================================================================
\section{Summary}
\label{sec:summary}

\begin{tcolorbox}[
  enhanced,
  colback=green!5!white,
  colframe=green!60!black,
  fonttitle=\bfseries,
  title={Lecture 6: Summary}
]
\begin{enumerate}
  \item \textbf{Without SUSY, three of four duality pillars are lost:}
    holomorphy, the exact beta function, and moduli space matching.
    Only 't~Hooft anomaly matching survives---exact, non-perturbative,
    and requiring no SUSY (\S\ref{sec:what-we-lose}).

  \item \textbf{Sannino--Schechter (1997):}  QCD degrees of freedom
    emerge from the auxiliary fields of the SUSY effective Lagrangian,
    \emph{suggesting} that duality structures may survive SUSY breaking
    (\S\ref{sec:sannino-schechter}).

  \item \textbf{Cheng--Shadmi tension (1998):}  Seiberg duality with
    large SUSY breaking produces tachyonic magnetic scalars---an
    \textbf{unresolved} tension.  Sannino's programme sidesteps this
    by conjecturing non-SUSY duality directly
    (\S\ref{ssec:cheng-shadmi}).

  \item \textbf{QCD dual (Sannino, 2009):}  Non-trivial anomaly
    matching solutions exist for non-SUSY QCD with a dual gauge
    group $\SU{2\Nf - 5N}$.  The conformal window prediction
    \emph{coincides} with the all-orders conjectured beta function
    (\S\ref{sec:qcd-dual}).  \emph{If the duality holds}, this
    provides a dual description of QCD.

  \item \textbf{SM-as-magnetic (Sannino, 2011):}  \emph{If the
    non-SUSY duality holds}, the SM is the magnetic dual of a
    purely fermionic electric theory with Pati--Salam structure.
    Consequences: no hierarchy problem, $n_g \geq 3$, common Yukawa
    origin.  Maekawa--Takahashi~\cite{MT1996} provided the
    direct SUSY precursor (\S\ref{sec:sm-magnetic}).

  \item \textbf{Emergent SUSY (Antipin--Mojaza--Pica--Sannino, 2011):}
    SUSY fixed points emerge from non-SUSY RG flow in the
    Veneziano limit.  This perturbative result \emph{suggests} that
    the quasi-SUSY spectrum is dynamically natural, but finite-$N$
    validity is uncontrolled (\S\ref{sec:emergent-susy}).

  \item \textbf{Central vulnerability:}  Non-SUSY duality has
    \textbf{no independent cross-validation beyond anomaly matching}.
    The construction postulates the duality and checks consistency.
    Consistency is not proof (\S\ref{sec:vulnerability}).
\end{enumerate}
\end{tcolorbox}


%% ========================================================================
%%  BIBLIOGRAPHY
%% ========================================================================
\begin{thebibliography}{99}

\bibitem{Seiberg1993}
  N.~Seiberg,
  ``Naturalness versus supersymmetric non-renormalization theorems,''
  \textit{Phys.~Lett.} \textbf{B318} (1993) 469
  [\texttt{hep-ph/9309335}].

\bibitem{NSVZ1983}
  V.~A.~Novikov, M.~A.~Shifman, A.~I.~Vainshtein, and
  V.~I.~Zakharov,
  ``Exact Gell-Mann--Low function of supersymmetric Yang--Mills
  theories from instanton calculus,''
  \textit{Nucl.~Phys.} \textbf{B229} (1983) 381.

\bibitem{tHooft1980}
  G.~'t~Hooft,
  ``Naturalness, chiral symmetry, and spontaneous chiral symmetry
  breaking,'' in \textit{Recent Developments in Gauge Theories},
  NATO ASI Series~B59 (Plenum, 1980) 135.

\bibitem{SS1997}
  F.~Sannino and J.~Schechter,
  ``Toy model for breaking super gauge theories at the effective
  Lagrangian level,''
  \textit{Phys.~Rev.} \textbf{D57} (1998) 170
  [\texttt{hep-ph/9708113}].

\bibitem{VY1982}
  G.~Veneziano and S.~Yankielowicz,
  ``An effective Lagrangian for the pure $\mathcal{N}=1$
  supersymmetric Yang--Mills theory,''
  \textit{Phys.~Lett.} \textbf{B113} (1982) 231.

\bibitem{CS1998}
  H.-C.~Cheng and Y.~Shadmi,
  ``Duality in the presence of supersymmetry breaking,''
  \textit{Nucl.~Phys.} \textbf{B531} (1998) 275
  [\texttt{hep-th/9801146}].

\bibitem{Sannino2009}
  F.~Sannino,
  ``QCD dual,''
  \textit{Phys.~Rev.} \textbf{D80} (2009) 065011
  [\texttt{arXiv:0907.1364}].

\bibitem{RS2008}
  T.~A.~Ryttov and F.~Sannino,
  ``Supersymmetry inspired QCD beta function,''
  \textit{Phys.~Rev.} \textbf{D78} (2008) 065001
  [\texttt{arXiv:0711.3745}].

\bibitem{Sannino2011}
  F.~Sannino,
  ``The Standard Model is natural as magnetic gauge theory,''
  [\texttt{arXiv:1102.5100}].

\bibitem{MT1996}
  N.~Maekawa and T.~Takahashi,
  ``Duality of SUSY model with Pati--Salam group,''
  \textit{Prog.~Theor.~Phys.} \textbf{95} (1996) 943
  [\texttt{hep-ph/9510426}].

\bibitem{PS1974}
  J.~C.~Pati and A.~Salam,
  ``Lepton number as the fourth `color',''
  \textit{Phys.~Rev.} \textbf{D10} (1974) 275.

\bibitem{AMPS2011}
  O.~Antipin, M.~Mojaza, C.~Pica, and F.~Sannino,
  ``Magnetic fixed points and emergent supersymmetry,''
  \textit{JHEP} \textbf{1306} (2013) 037
  [\texttt{arXiv:1105.1510}].

\bibitem{ASU2008}
  A.~Armoni, M.~Shifman, and M.~Unsal,
  ``Planar limit of orientifold field theories and emergent
  center symmetry,''
  \textit{Phys.~Rev.} \textbf{D77} (2008) 045012
  [\texttt{arXiv:0712.0672}].

\bibitem{BT2022}
  Y.~Bi and J.~Terning,
  ``Chiral gauge dynamics: candidates for non-supersymmetric
  dualities,''
  [\texttt{arXiv:2208.07842}].

\bibitem{CacciapagliaSannino2024}
  G.~Cacciapaglia and F.~Sannino,
  ``Charting SM duality and its signatures,''
  [\texttt{arXiv:2407.17281}].

\bibitem{Seiberg1995}
  N.~Seiberg,
  ``Electric--magnetic duality in supersymmetric non-Abelian gauge
  theories,''
  \textit{Nucl.~Phys.} \textbf{B435} (1995) 129
  [\texttt{hep-th/9411149}].

\end{thebibliography}

\end{document}
